\section{Statement and Proof}

We work in the following setting. Let $w\in\mathbb{R}_{>0}^m$ be a vector of positive weights with $\sum_{i=1}^m w_i=1$, and for a fixed $p\in(0,1)$ let $v\in\{0,1\}^m$ be a random vector whose coordinates are i.i.d.\ Bernoulli$(p)$. Set $X=\langle w,v\rangle$, so that $\mathbb{E}[X]=p$.

\begin{definition}
We say that \emph{Question~Q holds for $p$} if for every $m\ge 1$ and every $w\in\mathbb{R}_{>0}^m$ with $\sum_{i=1}^m w_i=1$, the random vector $v$ with i.i.d.\ Bernoulli$(p)$ coordinates satisfies
\[
\mathbb{P}\big[\langle v,w\rangle \ge p\big]\ge p.
\]
\end{definition}

\begin{conjecture}[Manickam--Mikl\'os--Singhi, parametrised form]\label{conj:MMS}
Let $c>0$. We say that the MMS conjecture \emph{holds for all $n\ge ck$} if: for all positive integers $n,k$ with $n\ge ck$ and any real numbers $x_1,\dots,x_n$ with $\sum_{i=1}^n x_i=0$, the number of $k$-element subsets $S\subseteq[n]$ with $\sum_{i\in S}x_i\ge 0$ is at least $\binom{n-1}{k-1}$.
\end{conjecture}

\begin{theorem}\label{thm:main}
Suppose there is a constant $c$ such that the MMS conjecture holds for all $n\ge ck$ (in the sense of Conjecture~\ref{conj:MMS}). Then Question~Q holds for all $p\in[0,1/c)$.
\end{theorem}

\begin{proof}
The cases $p=0$ is trivial (the inequality $\mathbb P[X\ge 0]\ge 0$ is automatic), so assume $p\in(0,1/c)$. We first treat rational $p$, then pass to the general case by a monotonicity/limit argument.

\medskip
\noindent\textbf{Step 1: Rational $p$.}
Suppose $p=k_0/n_0$ with $k_0,n_0$ positive integers and $p<1/c$. Fix an arbitrary $w\in\mathbb{R}_{>0}^m$ with $\sum_{i=1}^m w_i=1$, and write
\[
w_{\min}:=\min_{1\le j\le m} w_j>0.
\]
Let $t$ be a positive integer to be chosen large, and set
\[
k=k_0 t,\qquad n=n_0 t,\qquad\text{so that}\qquad \frac{k}{n}=\frac{k_0}{n_0}=p.
\]

\emph{Construction of an auxiliary matrix.}
Let $M$ be a random $n\times m$ matrix with entries in $\{0,1\}$, generated as follows. For each column $j\in[m]$ independently, choose a row index $R_j\in[n]$ uniformly at random and put a $1$ in position $(R_j,j)$ and $0$ in the other $n-1$ entries of column $j$. Thus each column has exactly one $1$. Denote the rows of $M$ by $r_1,\dots,r_n\in\{0,1\}^m$, and set
\[
y_i=\langle w,r_i\rangle\qquad (i\in[n]).
\]
Independently of $M$, choose a uniformly random $k$-element subset $S\subseteq[n]$, and define
\[
Y=\sum_{i\in S}y_i=\Big\langle w,\sum_{i\in S}r_i\Big\rangle.
\]

\medskip
\noindent\textbf{Claim 1.} \emph{$Y$ has the same distribution as $X=\langle w,v\rangle$.}

\begin{proof}[Proof of Claim 1]
Fix any deterministic $k$-set $S\subseteq[n]$, and consider the random vector $u^S:=\sum_{i\in S}r_i\in\{0,1\}^m$ (the randomness coming from $M$). For each column $j$, the $j$-th coordinate of $u^S$ equals $1$ if and only if the unique $1$ of column $j$ lies in one of the rows indexed by $S$, i.e.\ if and only if $R_j\in S$. Since $R_j$ is uniform on $[n]$, this happens with probability $|S|/n=k/n=p$. Moreover the columns are independent, so the coordinates of $u^S$ are independent Bernoulli$(p)$ variables. Hence, conditionally on $S$, the vector $u^S$ has exactly the distribution of $v$, and therefore $\langle w,u^S\rangle$ has the distribution of $X$. As this holds for every fixed $S$ and $S$ is independent of $M$, averaging over $S$ shows that $Y=\langle w,u^S\rangle$ has the same distribution as $X$.
\end{proof}

\medskip
\noindent\textbf{Claim 2.} \emph{If $t$ is large enough then $\mathbb{P}[Y\ge p]\ge p$.}

\begin{proof}[Proof of Claim 2]
Condition on the matrix $M$ (equivalently on the values $y_1,\dots,y_n$); we will show that the conditional probability $\mathbb P[Y\ge p\mid M]\ge p$ for every realisation of $M$, which then yields the unconditional bound.

Fix a realisation of $M$ and define
\[
z_i=y_i-\frac{1}{n}\qquad (i\in[n]).
\]
We record three properties.

\smallskip
\emph{(i) The $z_i$ sum to zero.} Since $M$ has exactly one $1$ in each of its $m$ columns, summing the rows gives $\sum_{i=1}^n r_i=\mathbf 1\in\{0,1\}^m$ (the all-ones vector). Hence
\[
\sum_{i=1}^n y_i=\Big\langle w,\sum_{i=1}^n r_i\Big\rangle=\langle w,\mathbf 1\rangle=\sum_{j=1}^m w_j=1,
\]
and therefore
\[
\sum_{i=1}^n z_i=\sum_{i=1}^n y_i-n\cdot\frac1n=1-1=0.
\]

\smallskip
\emph{(ii) Each $y_i$ is either $0$ or at least $w_{\min}$.} The entry $y_i=\langle w,r_i\rangle$ is a sum of those $w_j$ for which $M_{ij}=1$. If row $i$ is all zeros then $y_i=0$; otherwise $y_i$ contains at least one term $w_j\ge w_{\min}$, so $y_i\ge w_{\min}>0$.

\smallskip
\emph{(iii) For $t$ large, the only negative value among the $z_i$ is $-1/n$.} Choose $t$ large enough that
\[
\frac1n=\frac{1}{n_0 t}< w_{\min}.
\]
(This is possible since $w_{\min}>0$ is a fixed positive constant determined by $w$, while $1/(n_0 t)\to 0$ as $t\to\infty$.) Then by (ii), if $y_i>0$ we have $y_i\ge w_{\min}>1/n$, so $z_i=y_i-1/n>0$; and if $y_i=0$ then $z_i=-1/n<0$. Hence each $z_i$ is either $-1/n$ (when $y_i=0$) or strictly positive (when $y_i>0$). In particular all the negative $z_i$'s have the common value $-1/n$.

\smallskip
Now $z_1,\dots,z_n$ are real numbers summing to $0$ by (i), and $n\ge ck$ because $k/n=p<1/c$. By hypothesis the MMS conjecture holds for all $n\ge ck$, so the number of $k$-subsets $S\subseteq[n]$ with $\sum_{i\in S}z_i\ge 0$ is at least $\binom{n-1}{k-1}$. Therefore, for a uniformly random $k$-set $S$,
\[
\mathbb P\Big[\sum_{i\in S}z_i\ge 0\ \Big|\ M\Big]\ \ge\ \frac{\binom{n-1}{k-1}}{\binom{n}{k}}=\frac{k}{n}=p.
\]
Finally, $\sum_{i\in S}z_i\ge 0$ is equivalent to
\[
\sum_{i\in S}y_i\ge \frac{|S|}{n}=\frac{k}{n}=p,
\]
i.e.\ to $Y\ge p$. Hence $\mathbb P[Y\ge p\mid M]\ge p$ for every realisation of $M$, and taking expectation over $M$ gives $\mathbb P[Y\ge p]\ge p$.
\end{proof}

Combining Claims 1 and 2, for every $w$ as above and every rational $p\in(0,1/c)$,
\[
\mathbb P\big[\langle v,w\rangle\ge p\big]=\mathbb P[X\ge p]=\mathbb P[Y\ge p]\ge p.
\]
(Note that property (iii), and hence the whole argument, requires choosing $t$ large depending on $w$; but Claim~1 holds for every $t$, so the equality of distributions of $X$ and $Y$ is valid for the chosen large $t$, and the conclusion $\mathbb P[X\ge p]\ge p$ does not depend on $t$.) Thus Question~Q holds for every rational $p\in[0,1/c)$.

\medskip
\noindent\textbf{Step 2: Irrational $p$.}
Let $p\in(0,1/c)$ be arbitrary (in particular irrational), and fix $w\in\mathbb{R}_{>0}^m$ with $\sum_j w_j=1$. Let $v$ have i.i.d.\ Bernoulli$(p)$ coordinates. Choose a sequence of rationals $p_i\uparrow p$ with $0<p_i<p<1/c$.

For each $i$ we construct a coupling. Let $U_1,\dots,U_m$ be i.i.d.\ Uniform$[0,1]$. Define
\[
v_j=\mathbf 1[U_j\le p],\qquad v_j^{(i)}=\mathbf 1[U_j\le p_i]\qquad(j\in[m]).
\]
Then $v$ has i.i.d.\ Bernoulli$(p)$ coordinates, $v^{(i)}$ has i.i.d.\ Bernoulli$(p_i)$ coordinates, and since $p_i\le p$ we have $v_j^{(i)}\le v_j$ for all $j$. Because all weights $w_j$ are positive, this gives the pointwise (almost sure) domination
\[
\langle v,w\rangle\ \ge\ \langle v^{(i)},w\rangle.
\]
Hence
\[
\mathbb P\big[\langle v,w\rangle\ge p_i\big]\ \ge\ \mathbb P\big[\langle v^{(i)},w\rangle\ge p_i\big]\ \ge\ p_i,
\]
where the last inequality is the conclusion of Step~1 applied to the rational $p_i<1/c$ and the (Bernoulli$(p_i)$-distributed) vector $v^{(i)}$.

Now the events $A_i:=\{\langle v,w\rangle\ge p_i\}$ are nested decreasingly: since $p_i\uparrow p$, we have $p_i\le p_{i+1}$, so $A_{i+1}\subseteq A_i$. Their intersection is
\[
\bigcap_{i} A_i=\Big\{\langle v,w\rangle\ge p_i\ \text{for all }i\Big\}=\big\{\langle v,w\rangle\ge \sup_i p_i\big\}=\big\{\langle v,w\rangle\ge p\big\},
\]
using $\sup_i p_i=p$. By continuity of probability from above,
\[
\mathbb P\big[\langle v,w\rangle\ge p\big]=\mathbb P\Big[\bigcap_i A_i\Big]=\lim_{i\to\infty}\mathbb P[A_i]=\lim_{i\to\infty}\mathbb P\big[\langle v,w\rangle\ge p_i\big]\ \ge\ \lim_{i\to\infty}p_i=p.
\]

Since $w$ was arbitrary, Question~Q holds for $p$. This completes the proof for all $p\in[0,1/c)$.
\end{proof}

\begin{remark}
In Step~2, the identity $\bigcap_i A_i=\{\langle v,w\rangle\ge p\}$ holds because for a fixed realisation of $v$, the value $\langle v,w\rangle$ is a fixed real number, and a real number is $\ge p_i$ for every $i$ with $p_i\uparrow p$ if and only if it is $\ge p$. The argument also shows directly, without invoking irrationality, that the conclusion holds for \emph{all} $p\in[0,1/c)$ once it is known for all rationals in $[0,1/c)$; thus Steps~1 and~2 together cover every $p\in[0,1/c)$.
\end{remark}

\begin{remark}
The proof uses the MMS conjecture only for sequences $z_1,\dots,z_n$ in which all negative entries are equal (to $-1/n$), by property~(iii) above. Hence the same conclusion follows from the a priori weaker statement asserting the MMS bound for such sequences.
\end{remark}

\begin{proposition}\label{prop:counterexample}
For every $p\in(1/3,1)$ with $p\neq 1/2$, Question~Q fails.
\end{proposition}

\begin{proof}
Suppose first $1/3<p<1/2$. Take $m=3$ and $w=(1/3,1/3,1/3)$. Then $\langle v,w\rangle=\tfrac13(v_1+v_2+v_3)\ge p>1/3$ holds if and only if $v_1+v_2+v_3\ge 2$ (since the possible values of $\langle v,w\rangle$ are $0,\tfrac13,\tfrac23,1$, and we need a value $\ge p$ with $1/3<p<1/2$, i.e.\ a value $\ge 2/3$). Therefore
\[
\mathbb P[\langle v,w\rangle\ge p]=\binom{3}{2}p^2(1-p)+p^3=3p^2(1-p)+p^3=p\,(3p-2p^2).
\]
Hence
\[
p-\mathbb P[\langle v,w\rangle\ge p]=p\big(1-3p+2p^2\big)=p(1-2p)(1-p)>0
\]
for $p\in(1/3,1/2)$, since each factor $p,(1-2p),(1-p)$ is positive there. Thus $\mathbb P[\langle v,w\rangle\ge p]<p$.

Now suppose $1/2<p<1$. Take $m=2$ and $w=(1/2,1/2)$. Then $\langle v,w\rangle=\tfrac12(v_1+v_2)\ge p>1/2$ holds if and only if $v_1+v_2=2$ (the value must be $\ge p>1/2$, and the only attainable such value is $1$). Therefore
\[
\mathbb P[\langle v,w\rangle\ge p]=p^2<p,
\]
since $0<p<1$. In both cases Question~Q fails.
\end{proof}